\section*{Chapitre \ref{ch:pl:data-quality-and-lineage} --- Qualité des données et lignage}

\begin{solution}{exo:pl:data-quality-and-lineage:1}
Quatre~: 91\,687 à 91\,690. La capture du chapitre~2 attend que l'autre ligne les comble, puis enregistre l'étendue du trou et passe
à la suite~; la règle de séquence de ce chapitre signale la première ligne qui suit, afin que les consommateurs sachent que la journée est incomplète à cet endroit.
\end{solution}

\begin{solution}{exo:pl:data-quality-and-lineage:2}
Sur un instrument à grand tick, la plupart des transactions consécutives s'affichent au même prix, si bien que l'écart absolu médian des écarts récents vaut
zéro, et tout mouvement, même d'un seul tick, serait à une infinité d'unités robustes~: chaque changement de tick serait un pic. Le plancher (0.1~\%
ici) dit que rien de plus petit qu'un mouvement plausible n'est jamais une anomalie.
\end{solution}

\begin{solution}{exo:pl:data-quality-and-lineage:3}
Une erreur de décimale~: un facteur exactement dix. La médiane prospective vaut toujours 300\,000, car les transactions suivantes sont revenues à ce niveau, si bien que la
transaction s'écarte des deux médianes~; un véritable mouvement vers 3\,000\,000 resterait à ce niveau et ne s'écarterait que de la médiane rétrospective.
\end{solution}

\begin{solution}{exo:pl:data-quality-and-lineage:4}
Un indicateur peut se tromper (le saut), une ligne fausse pour un usage est nécessaire à un autre, et une suppression ne peut être ni
annulée. Les systèmes de surveillance et de reconstitution des transactions ont besoin des transactions erronées et annulées~: elles font partie de ce qui s'est passé.
\end{solution}

\begin{solution}{exo:pl:data-quality-and-lineage:5}
Erreurs conservées~: les erreurs de décimale qui multiplient un prix par dix dominent la moyenne pondérée par les volumes (quelques transactions à dix fois le prix). Avec
les lignes d'erreur exclues, le reste, $-0.02~\%$, provient des transactions perdues dans le trou de 200 événements~: l'échantillon est complet sauf là,
et aucun indicateur ne peut restituer des données manquantes.
\end{solution}

\begin{solution}{exo:pl:data-quality-and-lineage:6}
Les versions construites à partir de l'ancienne clôture~: les valorisations, le P\&L et les variables de recherche, qui utilisaient la clôture. Les barres à la minute, construites à partir
des seules transactions, n'ont pas besoin d'être reconstruites. Le lignage le sait, car chaque version enregistre les empreintes de contenu de ses entrées~: tout ce qui
se trouve en aval de la version de l'ancienne clôture est obsolète.
\end{solution}

\begin{solution}{exo:pl:data-quality-and-lineage:7}
Règle naïve~: 24 faux indicateurs de pic (les transactions après le saut de 3~\% de SIM2, jusqu'à ce que la médiane rétrospective rattrape) et un faux indicateur d'obsolescence
(la suspension de cotation de 60 secondes de SIM2). Règle informée~: aucun~; chaque défaut inséré est repéré par les deux règles.
\end{solution}

\begin{solution}{exo:pl:data-quality-and-lineage:8}
Elle supprime de véritables mouvements (un saut de 5~\% est une information sur une journée agitée, et chaque transaction qui le suit jusqu'à la suivante) et invente des cotations
que le marché n'a jamais affichées (un instrument suspendu ne cotait pas~; répéter la dernière cotation le fait paraître négociable). Les données brutes sont perdues,
de sorte qu'aucune des deux erreurs ne peut être retrouvée ensuite. Signaler plutôt, avec des règles qui connaissent l'état du marché et confirment les pics dans les deux sens, et conserver
les données telles que livrées.
\end{solution}

\begin{solution}{pb:pl:data-quality-and-lineage:1}
\begin{enumerate}
\item Complétude (trous de séquence, cotations périmées), validité (prix et tailles positifs), cohérence (carnet croisé, pics de prix, le
nombre de transactions du fournisseur), ponctualité (les délais de livraison du contrat~; la règle d'obsolescence en est la cousine intrajournalière).
\item L'accord du producteur et des consommateurs sur le schéma, les règles et les délais de livraison~: un changement devient un changement négocié, non une
surprise.
\item 47\,472 événements, 9\,258 cotations et 1\,190 transactions.
\item Les transactions annulées~: ce sont des prix valides de la journée~; seule la liste d'annulations de la plateforme, livrée plus tard, les identifie.
\item Ce n'est pas un échec des données du jour mais une nouvelle version d'un fait~: elle est stockée comme version, et le lignage trouve ce qu'il faut
reconstruire.
\item La médiane rétrospective retarde sur le nouveau niveau, si bien que chaque nouvelle transaction s'écarte d'environ 3~\% en unités robustes de quelques points de base~: 24 indicateurs
pour la graine~0.
\item 495 pour la règle naïve, 0 pour la règle informée.
\item Soixante secondes sans cotation dépassent 30 secondes~; la règle informée ne compte pas le temps de suspension de cotation.
\item Deux douzaines de transactions véritables après l'annonce ont été supprimées, et le mouvement de la journée a été sous-estimé.
\item Une règle qui repère chaque défaut inséré peut quand même signaler des centaines de lignes authentiques~; seuls les événements authentiques qui leur ressemblent permettent de mesurer les faux indicateurs.
\item Erreurs conservées~: SIM1 $+2.69~\%$, SIM2 $+1.08~\%$~; erreurs exclues~: à 0.03~\% près (SIM2 $-0.02~\%$).
\item Environ 24\,000.
\item 16 minutes signalées, celles où le trou a fait perdre des transactions~; il indique où regarder, non quelles transactions manquent.
\item Un indicateur marque des lignes et laisse la partition lisible~; une \omterm{def:pl:data-quality-and-lineage:flag}{quarantaine} retient la partition jusqu'à ce qu'une personne la libère, avec un
motif, ou qu'elle soit reconstruite.
\item Les valorisations, le P\&L et les variables.
\item \textbf{Résultat principal.} Sur vingt graines, les deux suites repèrent les 340 défauts insérés~; la suite naïve signale aussi 495 transactions authentiques
comme pics et 20 suspensions de cotation comme flux périmés, la suite informée n'en signale aucune. Les défauts étant conservés, les VWAP sont trop élevés de 2.69~\% et 1.08~\%, et la variance réalisée de SIM2
est 24\,000 fois trop grande~; avec les lignes d'erreur exclues, les VWAP sont à 0.03~\% près et la variance est exacte.
\item Confirmer les pics par rapport à la médiane prospective~: 495 des 515 faux indicateurs.
\item Sans la règle, 20\,000 transactions à 60~\% de leurs prix, pendant des années~; avec elle, les mêmes transactions sont signalées comme annulées, exclues de la
recherche et conservées pour la surveillance.
\item Pour chaque version de jeu de données, les empreintes de contenu des versions d'entrée dont elle est issue (et du code), de sorte qu'une nouvelle
version d'entrée désigne tout ce qui se trouve en aval.
\item Une couche de qualité des données consigne ce qu'elle trouve et laisse les personnes décider~; un script de nettoyage modifie les données et ne laisse aucune trace.
\end{enumerate}
\end{solution}

\begin{solution}{iq:pl:data-quality-and-lineage:1}
Trous de séquence, bornes de validité, cotations croisées ou bloquées, pics confirmés dans les deux sens, instruments périmés alors que le marché est ouvert,
transactions confrontées à la liste d'annulations de la plateforme, nombres et volumes comparés à une seconde source et à la clôture officielle~; chacun avec un
responsable, et des résultats sous forme d'indicateurs.
\iqlookfor{Les quatre familles, l'état du marché, des indicateurs et non des suppressions.}
\end{solution}

\begin{solution}{iq:pl:data-quality-and-lineage:2}
Une mauvaise transaction revient en général~: on la compare aux prix avant et après, en unités robustes avec un plancher~; un mouvement véritable persiste.
On vérifie aussi par rapport aux cotations, aux autres plateformes et aux nouvelles, et par rapport aux annulations de la plateforme publiées plus tard.
\iqlookfor{Biais d'anticipation dans un contrôle par lots~; le retournement comme signature.}
\end{solution}

\begin{solution}{iq:pl:data-quality-and-lineage:3}
La règle peut être fausse, d'autres consommateurs ont besoin des lignes, et une suppression ne peut être ni auditée ni annulée. Les indicateurs laissent chaque consommateur choisir et
permettent aux règles de s'améliorer.
\iqlookfor{Réversibilité et consommateurs différents.}
\end{solution}

\begin{solution}{iq:pl:data-quality-and-lineage:4}
Stocker la correction comme nouvelle version de la clôture, puis utiliser le lignage pour trouver chaque résultat construit à partir de l'ancienne version (valorisations, P\&L,
risque, variables) et les reconstruire, en signalant la variation du P\&L comme une correction, sans la remplacer en silence.
\iqlookfor{Versions et lignage~; la variation signalée.}
\end{solution}

\begin{solution}{iq:pl:data-quality-and-lineage:5}
Insérer des défauts connus et, tout aussi important, des événements authentiques qui leur ressemblent (suspensions de cotation, sauts, enchères, clôtures anticipées), et mesurer la détection
et les faux indicateurs par règle~; suivre les taux d'indicateurs dans le temps sur des données réelles et examiner chaque semaine un échantillon d'indicateurs.
\iqlookfor{Facteurs confondants et taux de faux indicateurs, et non seulement la détection.}
\end{solution}

\begin{solution}{iq:pl:data-quality-and-lineage:6}
Nommer chaque version de jeu de données par son empreinte de contenu~; enregistrer pour chacune ses versions d'entrée, sa version de code, ses paramètres et son exécution~; stocker les résultats avec
le manifeste~; répondre à «~de quoi ceci est-il construit~» et à «~qu'est-ce qui dépend de ceci~» par des requêtes sur le graphe~; recalculer à partir du manifeste après une
correction.
\iqlookfor{Versions adressées par leur contenu et un graphe interrogeable.}
\end{solution}
